Research on sequential decision-making under uncertainty contains two intellectual cultures. The \emph{inference culture} seeks to understand the world. A model in this culture specifies preferences, beliefs, constraints, and an equilibrium concept. The \emph{control culture} seeks to act in the world. Its operational model may instead consist of a transition kernel $P(s'|s,a)$ and a reward function $r(s,a)$. These mathematical objects serve different purposes.

The divide concerns the purpose of the model, not the profession, discipline, or immediate decision-maker. The inference culture primarily helps people understand the world and make decisions. The control culture primarily helps machines understand the world and make decisions. A common application pattern follows from this difference. Inference-oriented analysis is especially valuable when proposed changes are large, institutionally messy, hard to reverse, and require agreement among many affected decision-makers. Control-oriented analysis is especially useful in controlled experimental environments where actions can be tested repeatedly, failures can be rolled back, and implementation does not require prior agreement among many people. These settings favor each culture but do not define it. People can use control methods, and machine-assisted analysis can serve inference.

Inference-oriented analysis specifies primitives, asks which features data identify, and estimates unknown quantities. Identification, estimation, and solution are separate tasks. Even after the primitives have been estimated under maintained assumptions, the analyst must solve the resulting decision problem to compute a policy. Control-oriented analysis starts with an objective and a description of the system's dynamics, which may be supplied or learned. It then asks which policy can be computed and deployed under the relevant constraints on resources, timing, and robustness. Proportional-integral-derivative control (PID), linear-quadratic regulation (LQR), model predictive control (MPC), and reinforcement learning (RL) differ in the information and computation they require.

The data regime does not define either culture. Inference can use experiments or simulators. Control can learn from a fixed observational log. In observational settings, sorting, market clearing, and unobserved factors can make identification central. When intervention is unavailable, identification and support determine what can be learned from observed behavior. When controlled experimentation is available, exploration can create informative variation. Computational cost may then become central. A project can use both cultures. It may use inference to estimate or validate a model, then use control to choose actions within it.

The two cultures also share a vocabulary (``agent,'' ``learning,'' ``model,'' ``policy'') whose meanings diverge in ways that create persistent confusion. This chapter systematically translates between them.

\subsection{Core Distinctions}
\label{subsec:core_distinctions}

In RL, \emph{prediction} means estimating $V^\pi(s)$ or $Q^\pi(s,a)$ for a fixed policy $\pi$. \emph{Control} means finding a policy that maximizes expected return \citep{sutton2018}. \emph{On-policy} methods evaluate or improve the policy that generates the data. \emph{Off-policy} methods evaluate or improve a target policy $\pi$ using data generated by a different behavior policy $\mu$. Off-policy evaluation estimates the average return from deploying $\pi$, subject to identification and support conditions. It asks what would happen on average if $\pi$ were deployed. This differs from an individual structural counterfactual.\footnote{\citet{oberst2019counterfactual} condition on a realized trajectory and infer the posterior exogenous variables of a structural causal model before changing the actions. They show that the transition and reward distributions of a Markov decision process (MDP) alone do not identify these individual counterfactual trajectories.}

\emph{Online RL} can collect new experience by interacting with the environment. \emph{Offline RL}, also called batch RL, must learn from a fixed dataset collected previously under one or more behavior policies \citep{Levine2020}. These labels describe data availability, not response time. In this survey, \emph{real-time RL} means that action selection is subject to a deadline. A frozen policy can act in real time without learning online. An online learner can update its parameters more slowly than it selects actions.

An environment \emph{model} predicts how the environment responds to an action. It may return a distribution over next states and rewards or only a sample from that distribution. A \emph{model-based} method uses this model for planning. A \emph{model-free} method estimates values or policies from experience without using a learned or given environment model for planning \citep{sutton2018}. A model-free algorithm can still train in a fully specified simulator. A model-based method can instead learn its environment model from real observations. The distinction concerns how the algorithm uses information, not whether a scientific model of the domain exists.

The RL labels ``model-free'' and ``model-based'' do not correspond to the reduced-form and structural distinction in econometrics. They describe whether the algorithm plans through an environment model. They do not state whether the analyst imposes structural assumptions about preferences or equilibrium. ``Model-free'' also does not mean ``assumption-free.'' In the fully observed setting considered here, both variants operate within an MDP, which imposes the Markov property on the state. In the inference culture, a model may instead include agents, preferences, exclusion restrictions, and an equilibrium concept. A researcher can therefore specify a structural model and use a model-free RL algorithm as a computational solver, as the structural-estimation chapter demonstrates.

Training and execution need not be separate phases. During \emph{training}, observations change a value function, policy, reward model, or environment model. During \emph{execution}, the current policy selects actions. An offline-trained policy may be frozen before deployment. An online system may instead update its parameters while it continues to select actions. DiDi's deployed dispatcher, for example, selected assignments every two seconds and updated its value table every ten seconds \citep{SadeghiEshkevari2022DiDiScalable}. A hotel revenue-management field experiment updated its policy daily from completed booking episodes while continuing to recommend capacity allocations \citep{Chen2023hotelrl}. Both systems acted and learned in a live environment.

The word \emph{inference} has two meanings here. In machine learning systems, model inference means running a forward pass that produces outputs from inputs without updating the fitted parameters. In statistics, inference concerns properties of unknown quantities, including standard errors, confidence intervals, and hypothesis tests. This survey uses ``inference'' in the statistical sense and ``execution'' or ``deployment'' for applying a trained model. Established terms such as ``Bayesian inference'' and ``variational inference'' retain their conventional meanings. Most RL convergence and sample-complexity results describe the learning procedure. By themselves, they do not establish deployed performance. Claims about deployed performance also require the deployment environment and data-generating process to satisfy the conditions of those results. Table~\ref{tab:lifecycle_grid} summarizes these distinctions.

\newcommand{\tabitem}{{\small$\cdot$}~}
\begin{table}[H]
\centering
\caption{The reinforcement learning lifecycle grid. Rows show the source of experience. Columns distinguish learning or policy improvement from evaluating or using a fixed target policy.}
\label{tab:lifecycle_grid}
\begin{tabular}{|>{\raggedright\arraybackslash}p{2.8cm}|>{\raggedright\arraybackslash}p{5.5cm}|>{\raggedright\arraybackslash}p{5.5cm}|}
\hline
 & Component learning or policy improvement & Fixed-policy evaluation or action \\
\hline
Simulator &
  \tabitem Self-play or simulated experience \newline
  \tabitem Model learning from simulated transitions &
  \tabitem Fixed-policy evaluation \newline
  \tabitem Stress testing under simulated conditions \\
\hline
Historical data &
  \tabitem Offline policy learning from logged data \newline
  \tabitem Reward-model or dynamics-model estimation from logs &
  \tabitem Off-policy evaluation of a fixed target policy on logged data \\
\hline
Live environment &
  \tabitem Online value or policy updates \newline
  \tabitem Adaptive bandit learning &
  \tabitem Frozen-policy deployment \newline
  \tabitem Controlled evaluation of a fixed policy \\
\hline
\end{tabular}
\end{table}

The cells in Table~\ref{tab:lifecycle_grid} are not stages that every system must pass through. They separate the source of experience from how it is used. An offline method trains from historical data without interacting with the environment. Off-policy evaluation may estimate nuisance functions from historical data while holding the target policy fixed. A fixed policy evaluated in a live experiment collects new outcomes but does not improve itself during the experiment. An online controller may learn and act at the same time.

The Tmall pricing study shows one path through the grid \citep{Liu2019}. The researchers first trained deep Q-network (DQN) and deep deterministic policy gradient (DDPG) policies from historical pricing decisions generated by specialists or rules. They then evaluated candidate policies on later logged records. No suitable pricing simulator was available, so the next evaluation used live prices on Tmall. This workflow moved from historical preparation to live action. Other applied RL systems need not follow the same route.

\subsection{Overlapping Terminology}
\label{subsec:overlapping_terminology}

This vocabulary does not divide cleanly between academic fields. The same object can support inference when it is used to explain observations or quantify uncertainty, and control when it is used to select actions. The relevant distinction is therefore the purpose of the object in a particular analysis.

An \emph{agent} is the entity whose actions interact with an environment. It may be a person, a firm, an institution, or an algorithm. A control-oriented model uses the agent boundary to specify which action rule is being optimized. An inference-oriented model uses observed actions to learn about preferences, beliefs, constraints, or causal effects. In both cases, a \emph{policy} is a decision rule $\pi_t(a|H_t)$ based on the available history, with $\pi(a|s)$ as the stationary Markov special case. A government rule can instead be part of the environment, or it can be the policy under study, depending on the model boundary.

The \emph{environment} contains the variables outside the agent that determine what the agent observes and what follows from an action. An environment model may be known, estimated, or accessible only through samples. A \emph{reward} is the per-period criterion supplied to or learned by a control system, and it may itself be stochastic or unknown. A \emph{utility function} represents preferences. Inference-oriented work may estimate, calibrate, or test either object, while control-oriented work uses the resulting criterion to rank actions. Reward learning and inverse reinforcement learning combine both purposes.

The \emph{return}
\[
G_t=\sum_{k=0}^{\infty}\gamma^k R_{t+k+1}
\]
is the discounted sum of rewards after time $t$, when the sum is well defined. Under policy $\pi$, conditioning on $S_t=s$ gives $V^\pi(s)=\mathbb{E}_\pi[G_t\mid S_t=s]$. An \emph{episode} is an interaction sequence ending at a terminal condition or reset. A \emph{trajectory} is a realized interaction sequence containing the states or observations, actions, and rewards recorded by the model, and it may be finite, truncated, or continuing. A \emph{rollout} is a trajectory generated to evaluate or improve a policy. These terms describe data objects rather than cultures.

\emph{Learning} may refer to belief updating, statistical estimation, environment-model fitting, or policy improvement. Some value-based RL algorithms use stochastic approximation to approach a Bellman fixed point, but this is not a definition of learning in RL. The same mathematics also appears in inference-oriented models. \citet{MarcetSargent1989} analyze recursive least-squares learning through the stability of an associated ordinary differential equation. \citet{borkar2000} use the ODE method to establish convergence for asynchronous adaptive-critic and Q-learning algorithms in the average-cost setting. The shared tool is stochastic approximation and ODE stability, not a division between disciplines.

Information acquisition can likewise serve either culture. \emph{Active learning} selects observations or labels that are expected to improve an estimator. \emph{Exploration} selects actions partly for the information they provide about later decisions. Bayesian experimentation derives this tradeoff from a model of beliefs and payoffs, while methods such as the upper confidence bound (UCB) and Thompson sampling give alternative action rules and performance guarantees. The method becomes inference-oriented when the target is understanding an unknown quantity and control-oriented when the target is cumulative performance.

Several terms retain narrower technical meanings. \emph{Bootstrapping} in statistics refers to resampling observations \citep{Efron1979}, while in RL it refers to updating an estimate from another current estimate, as in the temporal-difference target $R_{t+1}+\gamma V(s_{t+1})$ \citep{sutton1988}. \emph{Function approximation} covers any parameterized representation of a value function, policy, reward, or model, including linear bases, kernels, and neural networks. \emph{Calibration} may mean selecting parameters to match empirical targets or checking whether predicted probabilities match observed frequencies. These meanings must be stated locally rather than assigned to one culture.

Bayesian and frequentist bandit formulations also cut across the two cultures. Thompson sampling allocates actions according to their posterior probability of being optimal \citep{Thompson1933}. In Bayesian dynamic allocation, posterior beliefs can form part of the state and the decision rule can account for the value of information \citep{Rothschild1974}. The Gittins index is optimal under the discounted, independent, rested-arm conditions for which its decomposition applies \citep{Gittins1979}. Frequentist analyses instead study guarantees such as stochastic-bandit pseudo-regret
\[
\bar{R}_T=\sum_{t=1}^{T}\mathbb{E}[\mu^*-\mu_{A_t}].
\]
Here $\mu_a=\mathbb{E}[R_t\mid A_t=a]$ and $\mu^*=\max_a\mu_a$ in the stationary stochastic-bandit model.
These are different formulations and criteria, not properties of different professions. Either can support inference about reward distributions or control of sequential allocation.

A \emph{contextual bandit} observes a context $x_t$, selects an action, and receives an immediate reward. The benchmark scores the action by this immediate reward rather than by a return that includes downstream state consequences. The context sequence may be independent, dependent over time, or adversarial, and an adaptive protocol may use past actions when choosing later contexts. Exogenous contexts are therefore an additional modeling assumption, not part of the definition. Under the immediate-reward benchmark, the algorithm receives no credit or penalty for downstream effects. If its actions change later contexts, the benchmark omits those consequences rather than making them absent.

\emph{Policy evaluation} can represent the same deployment-value functional at different levels of conditioning. In an inference-oriented analysis, it often estimates the average outcome under a specified decision rule, using experimental or observational data under explicit identification and support conditions. In a control-oriented analysis, it computes or estimates the state-conditional value $V^\pi(s)$ for a fixed policy from a model or experience. Averaging $V^\pi(s)$ over the deployment initial-state distribution yields the average policy value when the definitions align. Neither formulation requires a simulator, and neither restricts the policy to a government rule or an algorithmic controller.

\emph{Identification} asks which model features the observable distribution determines. A feature is point identified when any two admissible models that generate the same distribution of observables assign it the same value \citep{Lewbel2019}. Set, local, generic, and weak identification relax or qualify this uniqueness requirement. The issue is central whenever an analysis interprets preferences, rewards, dynamics, or causal effects. System identification estimates transition dynamics for later control \citep{RossBagnell2012}. Inverse reinforcement learning makes the same issue explicit because different rewards can induce observationally equivalent behavior \citep{KimGarg2021}. The companion survey develops this connection to dynamic discrete choice \citep{RustRawat2026}.

\emph{Regret} and \emph{efficiency} also require an explicit comparator. Regret may compare cumulative reward with the best fixed action, the best policy in a class, or an optimal MDP policy. Statistical efficiency concerns sampling variance or information bounds. Sample efficiency concerns the number of observations or interactions needed to reach a performance criterion, and computational efficiency concerns time or memory. Welfare efficiency is a separate property of an allocation. None of these meanings belongs exclusively to either culture.

The \emph{discount factor} is both part of a sequential objective and a mathematical device. In preference models it may be estimated, calibrated, or maintained as an assumption. Exponential discounting in \citet{Koopmans1960} follows from an axiom system combining ordering, continuity, sensitivity, stationarity, constant equivalence, initial-tradeoff independence, and tail conditions, not from impatience alone \citep{Bleichrodt2008}. In discounted RL, $\gamma$ sets the weight on delayed rewards and the effective horizon, while $\gamma<1$ makes the discounted Bellman operator a contraction on bounded value functions under the sup norm. \citet{Pitis2019} derives a more general state-action-dependent discount representation and shows that an optimizing MDP with fixed $\gamma<1$ can match optimal utility and the optimal policy, while preferences over suboptimal policies can differ.

\subsection{Structural Equivalences}
\label{subsec:structural_equivalences}

The softmax map appears in both cultures under explicit assumptions. Let $q(s,a)$ be a finite collection of action scores. The map
\begin{equation}
\label{eq:softmax_logit}
\pi_q(a \mid s) = \frac{\exp(q(s,a) / \tau)}{\sum_{a' \in \mathcal{A}} \exp(q(s,a') / \tau)},
\end{equation}
with $\tau>0$ is also the multinomial-logit probability when $q(s,a)$ is systematic utility and the additive shocks are independent Type I extreme value draws with common scale $\tau$ and a common location normalization \citep{McFadden1974}. The scores may instead be policy logits, ordinary action-value estimates, or entropy-regularized optimal action values. The formula alone does not determine which interpretation applies. As $\tau\to0$, probability concentrates on the maximizing actions and becomes deterministic only when the maximizer is unique.

The exact common structure is the variational identity
\begin{equation}
\label{eq:soft_value}
\mathcal{L}_{\tau}(q)
=\tau\log\sum_{a\in\mathcal{A}}\exp\!\left(\frac{q_a}{\tau}\right)
=\max_{p\in\Delta(\mathcal{A})}
\left\{\sum_a p_a q_a+\tau H(p)\right\},
\end{equation}
where $H(p)=-\sum_a p_a\log p_a$ and the unique optimizer is the softmax map in Equation~\eqref{eq:softmax_logit}. Positive entropy appears as a bonus in the primal problem, while negative entropy is the convex conjugate penalty of log-sum-exp. In a dynamic discrete choice model with systematic choice-specific values $v(x,d)$ and location-zero Gumbel shocks of scale $\tau$, the expected maximum is
\[
\mathbb{E}\!\left[\max_d\{v(x,d)+\varepsilon_d\}\right]
=\mathcal{L}_{\tau}(v(x,\cdot))+\tau\gamma_{\mathrm E},
\]
where $\gamma_{\mathrm E}$ is Euler's constant. Centering the shocks at mean zero removes this additive constant, as in the normalization used by \citet{Rust1987}.

The dynamic discrete choice (DDC) choice-specific value assumes optimal continuation after the current choice. Under the matched entropy convention, its nearest RL analogue is the optimal soft action value $Q^*_{\tau}$, not $Q^\pi$ for an arbitrary fixed policy. Under the logit model,
\begin{equation}
\label{eq:log_ccp_normalization}
v(x,d)-\mathcal{L}_{\tau}(v(x,\cdot))
=\tau\log P(d\mid x),
\end{equation}
and hence
\[
v(x,d)-v(x,d_0)
=\tau\left[\log P(d\mid x)-\log P(d_0\mid x)\right].
\]
This is the Hotz-Miller inversion for choice-specific value differences \citep{HotzMiller1993}. It is a soft normalization by the inclusive value, not the ordinary advantage $Q^\pi-V^\pi$.

The convex duality is exact after its scale and domain are stated. The conjugate of $\mathcal{L}_{\tau}$ is $\tau\sum_a p_a\log p_a$ on the probability simplex and is infinite outside it. \citet{ChiongGalichonShum2016} use this duality to connect choice probabilities with normalized choice-specific payoffs, and \citet{FosgerauSorensen2022} extend the inversion beyond logit. The forward control problem and the inverse estimation problem remain different tasks with different inputs.

The same softmax map can express different purposes. \citet{WilliamsPeng1991} study an entropy-maximizing REINFORCE variant. Maximum causal entropy and control as inference give probabilistic or variational interpretations of regularized control \citep{Ziebart2010,Levine2018}. A controller may randomize because entropy is part of its objective, even when the environment is known. In a random-utility model, choice is deterministic conditional on the realized utility shocks, while the analyst integrates over shocks that are unobserved. This is a distinction between a control objective and an inference model, not between professions.

\subsection{Notation}
\label{subsec:notation}

Table~\ref{tab:notation} records nearby notation without treating the objects as universal equivalents. The control column describes objects used to select or evaluate actions. The inference column describes objects used to interpret behavior or estimate unknown quantities.

\begin{table}[H]
\centering
\caption{Conditional notation map between control and inference uses. The final column states the condition under which the objects align.}
\label{tab:notation}
\small
\begin{tabular}{>{\raggedright\arraybackslash}p{2.6cm}>{\raggedright\arraybackslash}p{2.9cm}>{\raggedright\arraybackslash}p{3.0cm}>{\raggedright\arraybackslash}p{5.0cm}}
\toprule
Object & Control notation & Inference notation & Relation \\
\midrule
State & $s_t\in\mathcal{S}$ & $x_t\in\mathcal{X}$ & A Markov or sufficient state, not an arbitrary covariate. \\
Action & $a_t\in\mathcal{A}$ & $d_t$ & Each denotes the current decision or control. \\
Flow criterion & $r(s,a)$ & $u_\theta(x,d)$ & Flow reward, payoff, or systematic utility. \\
Discount factor & $\gamma$ & $\beta$ & Weights on delayed flow; endpoints depend on horizon and objective. \\
Policy or rule & $\pi_t(a\mid H_t)$ & $g_\theta(x,\varepsilon)$ & Kernel versus rule conditional on private shocks. \\
Action probability & $\pi(a\mid s)$ & $P_\theta(d\mid x)$ & The probabilities coincide only under a specified stochastic-policy or random-utility model. \\
Fixed-rule value & $V^\pi(s)$ & $V^g_\theta(x)$ & Expected return under specified rule $g$. \\
Optimal soft value & $V^*_{\tau}(s)$ & $\bar V_\theta(x)$ & Ex-ante optimum when primitives and entropy or shock normalization match. \\
Choice-specific value & $Q^*_{\tau}(s,a)$ & $v_\theta(x,d)$ & Current flow plus optimal continuation under matched entropy and logit models. \\
Realized return & $G_t$ & $\sum_{k\geq0}\beta^k u_{t+k}$ & Discounted sum under the stated timing convention. \\
Transition kernel & $P(s'\mid s,a)$ & $f_\theta(x'\mid x,d)$ & Conditional laws of the next state. \\
\bottomrule
\end{tabular}
\end{table}

Discount endpoints depend on the horizon and objective. A myopic model can use zero, a finite terminating problem can use an undiscounted sum, and a discounted infinite-horizon contraction requires a factor below one.

The temporal-difference (TD) innovation is distinct from the population Bellman residual. For policy evaluation, let $\delta_t=R_{t+1}+\gamma V(S_{t+1})-V(S_t)$ and $(\mathcal{T}^\pi V)(s)=\mathbb{E}_\pi[R_{t+1}+\gamma V(S_{t+1})\mid S_t=s]$. When actions and transitions are sampled under $\pi$,
\[
\mathbb{E}[\delta_t\mid S_t=s]
=(\mathcal{T}^\pi V)(s)-V(s).
\]
Mean squared sampled TD error therefore includes the target's conditional variance and is not generally the mean squared Bellman error. This produces the double-sampling problem in direct Bellman-error minimization \citep{sutton2018}.
